\section*{第 \ref{ch:g6:decimals} 章 --- 小数}

\begin{solution}{exo:g6:decimals:1}
$5.37$; \quad $0.9$; \quad $12.004$; \quad $8.05$.
\end{solution}

\begin{solution}{exo:g6:decimals:2}
十分位上的数字是 $3$；百分位上的数字是 $5$；$8$ 位于十位上。
\[
86.354 = 8 \times 10 + 6 + 3 \times \frac{1}{10} + 5 \times
\frac{1}{100} + 4 \times \frac{1}{1000}.
\]
\end{solution}

\begin{solution}{exo:g6:decimals:3}
$5.3 > 5.29$（比较 $5.30$ 与 $5.29$）；\quad
$0.7 = 0.70$（末尾的零不改变大小）；\quad
$2.09 < 2.9$（十分位：$0 < 9$）；\quad
$14.5 > 14.49$。
\end{solution}

\begin{solution}{exo:g6:decimals:4}
补写成两位小数：$4.20$；$4.05$；$4.51$；$4.15$；$4.50$。排序如下：
\[
4.05 < 4.15 < 4.2 < 4.5 < 4.51 .
\]
\end{solution}

\begin{solution}{exo:g6:decimals:5}
每个刻度表示十分之一。$A$：$6.3$；\quad $B$：$6.7$；\quad
$C$：$7.2$。
\end{solution}

\begin{solution}{exo:g6:decimals:6}
$45.80 + 7.65 = 53.45$.

$23.40 - 8.72 = 14.68$（验算：$14.68 + 8.72 = 23.4$）。

$56 \times 24 = 1344$，两个因数的小数位数共两位：
$5.6 \times 2.4 = 13.44$.
\end{solution}

\begin{solution}{exo:g6:decimals:7}
$6.42 \times 10 = 64.2$; \quad
$0.35 \times 1000 = 350$; \quad
$78.1 \div 100 = 0.781$; \quad
$5 \div 1000 = 0.005$.
\end{solution}

\begin{solution}{exo:g6:decimals:8}
$2.4$ m $= 240$ cm；\quad $370$ g $= 0.37$ kg（除以 $1000$）；
\quad $0.85$ km $= 850$ m；\quad $56$ mm $= 0.056$ m。
\end{solution}

\begin{solution}{exo:g6:decimals:9}
精确到个位：$24$（下一位是 $8$，应进一）。精确到十分位：$23.9$。
精确到百分位：$23.87$。$9.97$ 精确到十分位时，下一位是
$7$，所以十分位上的 $9$ 进一 --- $10.0$。
\end{solution}

\begin{solution}{exo:g6:decimals:10}
法棍的价格：$3 \times 1.15 = 3.45$。找零：
$5 - 3.45 = 1.55$。
\end{solution}

\begin{solution}{exo:g6:decimals:11}
在 $7.4$ 与 $7.5$ 之间，例如 $7.45$（或 $7.41$、$7.499$，
\dots）。在 $3.99$ 与 $4$ 之间，例如 $3.995$。两种情况下
这样的数都有 \emph{无穷多个}：总能再多写
几位小数。
\end{solution}

\begin{solution}{exo:g6:decimals:12}
最大的数：数字按从大到小排列，小数点尽可能靠后，
即 $85.2$。最小的数：数字按从小到大排列，小数点尽可能靠前，
即 $2.58$。
\end{solution}

\begin{solution}{pb:g6:decimals:1}
\textbf{1.} 补零后可见 $3$ 较大：$3.000 > 2.999$。两者
之差为 $3.000 - 2.999 = 0.001$，即千分之一。

\textbf{2.} 严格介于 $7.4 = 7.40$ 与 $7.5 = 7.50$ 之间的数有
\[
7.41,\ 7.42,\ 7.43,\ 7.44,\ 7.45,\ 7.46,\ 7.47,\ 7.48,\ 7.49 :
\]
共九个数，每个新刻度上各一个。

\textbf{3.} 图景相同，只是缩小为原来的十分之一：在 $7.440$
与 $7.450$ 之间的数有 $7.441, 7.442, \dots, 7.449$ ---
仍是九个数。

\textbf{4.} 在 $7.400$ 与 $7.500$ 之间，有三位小数的数为
$7.401, 7.402, \dots, 7.499$：全部
千分数，即 $401$ 至 $499$ 的全部，共 $99$ 个数。

\textbf{5.} “放大”的方法：把两个数写成相同的
小数位数，然后 \emph{再添加一位小数}
--- 在较小的数后面添上 $1$ 到 $9$ 中的一个数字，所得的数
严格介于两者之间。事实上 $5.67 = 5.670$ 与
$5.68 = 5.680$，且
\[
5.670 < 5.675 < 5.680 ;
\]
同理 $0.1999 = 0.19990 < 0.19995 < 0.20000 = 0.2$。相邻两个
刻度之间，总有一整级
由九个更细的刻度构成。

\textbf{6.} $3 < 3.05 < 3.1$；接着 $3 < 3.005 < 3.01$；再
$3 < 3.0005 < 3.001$。每个候选数都会在再一次“放大”后被排除。

\textbf{7.} 无论汤姆提出什么数——称之为他的
候选数，它严格大于 $3$——由第 5 题的方法
可得一个严格介于 $3$ 与该候选数之间的数。因此
该候选数\emph{并非}离 $3$ 最近的数：还存在更接近的数。
既然对每一个候选数都无例外地如此，
便不存在任何数可以是“紧接在 $3$ 之后的数”：
这样的数根本不存在。

\textbf{8.} $2.9 < 2.95 < 3$，\ $2.99 < 2.995 < 3$，\
$2.999 < 2.9995 < 3$。且 $3 - 2.999 = 0.001$。若写二十个
9，则差为小数点之后接十九个零再跟一个
$1$——即位于小数点后第二十位的一个单位：
极其微小，但并不为零。无论汤姆写多少个 9，
他都永远到不了 $3$。

\textbf{9.} 整数以 $1$ 为步长递增：大于 $7$
而小于 $8$，则意味着存在整数个单位严格位于
$7$ 与 $8$ 之间，而这样的整数并不存在。“放大”之所以能避开这一点，
仅仅在于写出\emph{小数点之后}的数字——
而这恰恰是整数所没有的。

\textbf{10.} 四舍五入后得到 $3.7$ 的最小长度是 $3.65$ cm：
其后一位数字是 $5$，按规则应向\emph{上}舍入
（见\cref{def:g6:decimals:rounding}）。同理，$3.75$ 向上舍入为 $3.8$。
因此，“精确到十分位的 $3.7$ cm”意味着：
实际长度至少为 $3.65$ cm，且严格小于 $3.75$ cm——
放大后的可能真实长度区间，隐藏在一个
四舍五入的近似值之后。

\textbf{11.} $0.0999 < 0.123 < 0.58 < 0.6$。一句话概括：
从左往右逐位比较，$0.0999$ 的十分位是 $0$，
而 $0.6$ 的十分位是 $6$，比较在此即已分出大小——
所写数字的个数并不能说明大小
（见\cref{met:g6:decimals:compare}）。

\textbf{12.} 不能。以分计，$4.99$ 欧元为 $499$ 分，$5$
欧元为 $500$ 分——相邻的\emph{整数}，
它们之间别无他数。价格恰好有两位小数，
实质上是伪装成小数的整数（以分为单位）：与第 9 题中的
整数一样，它们确实有紧挨着的“邻居”。无穷的“放大”
需要的是无限增加小数位数的权利。

\textbf{13.} 接近 $6$ 的候选数是 $5.82$ 和 $8.25$（一个
以 $2$ 开头、或以 $58$、$25$、$82$、$85$ 开头的数，离 $6$ 都很远）。
距离分别为 $6 - 5.82 = 0.18$ 与 $8.25 - 6 = 2.25$。
最接近的是 $5.82$。

\textbf{14.} 例如 $3.1415$，因为
$3.1410 < 3.1415 < 3.1420$；再如 $3.14159$，因为
$3.14150 < 3.14159 < 3.14160$。（这两者都是
在真正逼近 $\pi = 3.14159\dots$ 的过程中的实际步骤）

\textbf{15.} 假设有人给出一个数，无论它有多大，
一次“放大”会使每一对相邻刻度
在 $0$ 与 $1$ 之间变成九个新的数，且这种“放大”可以
无限重复——第 2、3、4 题只是前两级。重复足够多次，
便会产生比所给数更多的
位于 $0$ 与 $1$ 之间的小数。因此，没有任何数大到
足以将它们一一数尽：这些小数有无穷多个。
\end{solution}
